% eksperymentalne tłumaczenie na włoski

%

\section{Problem Hansena} % lang-pl
\index{problem!Hansena}% % lang-pl
W trygonometrii problem Hansena to zagadnienie z zakresu pomiarów w terenie, nazwane na cześć astronoma Petra Andreasa Hansena, który będzie pracował przy mapach państwa duńskiego. % lang-pl
\section{Problema di Hansen} % lang-it
\index{problema!di Hansen}% % lang-it
In trigonometria, il problema di Hansen è un problema di topografia, intitolato all'astronomo Peter Andreas Hansen, che lavorò alla realizzazione delle carte del Regno di Danimarca. % lang-it

\begin{problem}
    Dane są dwa znane punkty $A$ i $B$ oraz dwa nieznane punkty $P_1$ i $P_2$. % lang-pl
    Z punktów $P_1$ i $P_2$ obserwator mierzy kąty utworzone przez linie łączące go z każdym z pozostałych trzech punktów skierowane do każdego z pozostałych trzech punktów. % lang-pl
    Wyznaczyć położenie punktów $P_1$ i $P_2$. % lang-pl
    Sono dati due punti noti $A$ e $B$ e due punti ignoti $P_1$ e $P_2$. % lang-it
    Dai punti $P_1$ e $P_2$ un osservatore misura gli angoli formati dalle rette che lo collegano a ciascuno degli altri tre punti. % lang-it
    Determinare la posizione dei punti $P_1$ e $P_2$. % lang-it
\end{problem}

Znamy cztery kąty $\alpha_1 = \angle AP_1B$, $\beta_1 = \angle BP_1 P_2$, $\alpha_2 = \angle  AP_2B$, $\beta_2 = \angle AP_2P_1$. % lang-pl
Wtedy jedno z~rozwiązań zaczyna się od wyznaczenia kątów % lang-pl
Sono noti i quattro angoli $\alpha_1 = \angle AP_1B$, $\beta_1 = \angle BP_1 P_2$, $\alpha_2 = \angle AP_2B$, $\beta_2 = \angle AP_2P_1$. % lang-it
Una delle soluzioni inizia determinando gli angoli % lang-it
\begin{equation}
    \gamma = \pi - \alpha_1 - \beta_1 - \beta_2, \quad
    \delta = \pi - \alpha_2 - \beta_1 - \beta_2
\end{equation}
oraz ilorazu % lang-pl
e il rapporto % lang-it
\begin{equation}
    k = \frac{\sin \gamma \sin \alpha_2 \sin \beta_1}{\sin \delta \sin \alpha_1 \sin \beta_2}.
\end{equation}
Niech % lang-pl
Sia % lang-it
\begin{equation}
    \theta = 2 \arctan \left(\frac{k-1}{k+1} \tan \frac{\beta_1 + \beta_2}{2} \right).
\end{equation}
Wprowadzamy pomocnicze kąty $\varphi = \frac 1 2 (\beta_1 + \beta_2 + \theta)$ oraz $\psi = \frac 1 2 (\beta_1 + \beta_2 - \theta)$. % lang-pl
Szukany odcinek ma długość % lang-pl
Introduciamo gli angoli ausiliari $\varphi = \frac 1 2 (\beta_1 + \beta_2 + \theta)$ e $\psi = \frac 1 2 (\beta_1 + \beta_2 - \theta)$. % lang-it
Il segmento cercato ha lunghezza % lang-it
\begin{equation}
    |P_1 P_2| = |AB| \cdot \frac{\sin \psi \sin \gamma}{\sin \alpha_1 \sin \beta_2} = |AB| \cdot \frac{\sin \varphi \sin \delta}{\sin \alpha_2 \sin \beta_1},
\end{equation}
przy czym by uniknąć błędów numerycznych bierzemy ten ułamek, którego mianownik ma większą wartość bezwzględną. % lang-pl
per evitare errori numerici si utilizza la frazione il cui denominatore ha valore assoluto maggiore. % lang-it

%